% ---------------------------------------------------------------------------- %
% Abstract
\chapter*{Abstract}

This work presents the development of a Self-Service Engineering Platform aimed at abstracting the complexity of provisioning and configuring infrastructure environments for research and development. Traditionally, creating virtual machines (VMs) or containers, as well as configuring specific software and services, demands time and specialized technical knowledge from the end user \cite{KunalBhandari}. The developed solution, named ProxLab, turns this manual process into an on-demand model: through a web portal, the user (student or researcher) specifies their infrastructure needs, and the backend orchestrates the creation of virtual machines on Proxmox VE through Terraform (Infrastructure as Code) or the provisioning of application containers on a dedicated Kubernetes (k3s) cluster, with per-profile quotas and access via web terminal or network. This Platform Engineering approach \cite{Sandobalín} empowers researchers who do not have in-depth infrastructure knowledge to provision their own work environments in an agile way, reducing dependency on operations teams and the occurrence of manual errors \cite{Gokhale}. The results obtained demonstrate the technical feasibility of the approach, evidencing gains in autonomy and reduction in provisioning time compared to the manual process \cite{Muhammad}.

\vfill \textbf{Keywords:} Platform Engineering, Self-Service, Infrastructure as Code (IaC), Terraform, Kubernetes, Proxmox, Infrastructure Automation.
\\

% ---------------------------------------------------------------------------- %\\